% Generated draft: requires mathematical and source review.
\documentclass[11pt]{article}
\usepackage[margin=25mm]{geometry}
\usepackage{fontspec}
\setmainfont{Latin Modern Roman}
\usepackage{amsmath,amssymb,mathtools,unicode-math}
\setmathfont{Latin Modern Math}
\usepackage{enumitem,needspace,xurl,hyperref}
\setlength{\parindent}{0pt}
\setlength{\parskip}{5pt}
\setlength{\emergencystretch}{3em}
\newcommand{\N}{\mathbb{N}}
\newcommand{\Z}{\mathbb{Z}}
\newcommand{\Q}{\mathbb{Q}}
\newcommand{\R}{\mathbb{R}}
\newcommand{\C}{\mathbb{C}}
\newcommand{\F}{\mathbb{F}}
\newcommand{\A}{\mathbb{A}}
\newcommand{\PP}{\mathbb{P}}
\newcommand{\E}{\mathbb{E}}
\newcommand{\eps}{\varepsilon}
\newcommand{\dd}{\,\mathrm{d}}
\begin{document}
\section*{UnsolvedMath ID: 30006646}
\textbf{Source problem code:} catalogue entry
\subsection{Definitions and mathematical statement}
Let \(k\) be a field. The supplied excerpt describes a finite tensor category over \(k\) as a \(k\)-linear tensor category \(\mathcal C\) satisfying the following conditions: \(\mathcal C\) is abelian and rigid; \(\dim_k\operatorname{Hom}_{\mathcal C}(X,Y)<\infty\) for every pair of objects \(X,Y\); every object has finite length; there are finitely many isomorphism classes of simple objects; and the tensor unit \(\mathbf 1\) is simple. Rigidity means that every object has left and right duals. The excerpt states no restriction on the characteristic or algebraic closure of \(k\).

For objects \(M,N\in\mathcal C\), set
\[
\operatorname{Ext}^{*}_{\mathcal C}(M,N)
=\bigoplus_{n\geq 0}\operatorname{Ext}^{n}_{\mathcal C}(M,N),
\qquad
H^{*}(\mathcal C)
=\operatorname{Ext}^{*}_{\mathcal C}(\mathbf 1,\mathbf 1),
\]
where higher extension groups may be understood as Yoneda extension groups, and the multiplication on \(H^{*}(\mathcal C)\) is the Yoneda product. One conventional left action of \(H^{*}(\mathcal C)\) on \(\operatorname{Ext}^{*}_{\mathcal C}(M,N)\) is
\[
h\cdot x=(h\otimes\operatorname{id}_N)\circ x,
\qquad
h\in H^p(\mathcal C),\quad
x\in\operatorname{Ext}^{q}_{\mathcal C}(M,N),
\]
using \(\mathbf 1\otimes N\cong N\) and Yoneda composition. Tensoring is exact in this abelian rigid setting. This action convention is made explicit here; the supplied excerpt does not specify a tensor side.

\textbf{Conjecture 1.2 (attributed in the source to Etingof--Ostrik, 2004).} For every finite tensor category \(\mathcal C\),
\begin{enumerate}
\item the cohomology algebra \(H^{*}(\mathcal C)\) is Noetherian;
\item for every pair of objects \(M,N\in\mathcal C\), the graded module \(\operatorname{Ext}^{*}_{\mathcal C}(M,N)\) is finitely generated over \(H^{*}(\mathcal C)\).
\end{enumerate}
Here finite generation means that there are finitely many homogeneous extension classes whose \(H^{*}(\mathcal C)\)-linear combinations exhaust the module. The source calls the conjunction of these properties \emph{cohomological finite generation}, abbreviated CFG.

\textbf{Unresolved formulation qualifications.} The displayed definition in the supplied excerpt omits an enough-projectives condition that is normally part of the finite abelian category hypothesis used for finite tensor categories. It also contains an unexplained symbol \(A\) in the finite-simple-objects condition, interpreted above as referring to \(\mathcal C\). Whether additional hypotheses were intended must be checked against the complete definition; they have not been silently added. The excerpt also does not specify whether Noetherianity refers to all ideals or only homogeneous ideals, or explicitly specify left/right conventions. These conventions and the module action require confirmation for a definitive formulation.

The source motivates CFG by its role in tensor triangular geometry for \(\mathcal C\) and its ind-completion, and mentions a further conjecture concerning \(\operatorname{stab}\mathcal C\). No statement of that further conjecture is supplied, so it is not included here.
\subsection{Short English statement}
The Etingof--Ostrik conjecture asserts that every finite tensor category has Noetherian cohomology and that the extensions between every pair of objects form a finitely generated module over that cohomology. The supplied definition and algebra-action conventions need confirmation. The record labels the problem partially solved; its dated assessment is supplied curation, not an independently verified mathematical conclusion.
\subsection{Sources}
\begin{itemize}
\item Local UnsolvedMath registry, supplied complete record: permanent ID \(30006646\), title \emph{Etingof--Ostrik Cohomological Finite-Generation Conjecture}, exact source code \texttt{catalogue entry}. The record belongs to the Oberwolfach Reports Open Problems collection and attributes Conjecture 1.2 to Etingof--Ostrik, 2004; no complete bibliographic citation for that original attribution is supplied.
\item \emph{Cohomology of Finite Groups: Interactions and Applications} (2026), Oberwolfach Reports, pp.\ 541--610, DOI \texttt{10.4171/owr/2026/9}: \url{https://doi.org/10.4171/owr/2026/9}. This is the supplied originating report. The record labels the statement \texttt{corrected\_verified} and says it was reconstructed and checked against the official report; that verification is reported from the record and has not been independently performed here.
\item Bergh, \emph{Support varieties without the tensor product property}: \url{https://doi.org/10.1112/blms.13048}. Listed among the supplied literature sources.
\item \emph{Separable equivalences, finitely generated cohomology and finite tensor categories}: \url{https://doi.org/10.1007/s00209-023-03309-3}. Listed among the supplied literature sources; the record supplies no author attribution for this item.
\item Literature assessment dated exactly \texttt{2026-08-22}: the supplied curation labels the problem partially solved and states that the conjecture remains open for arbitrary finite tensor categories, while being proved in important classes including braided categories and several Hopf-algebraic cases. This is supplied dated curation, not independently verified here; the scope and hypotheses of these claimed special cases require independent checking.
\item The excerpt refers to its definition as Definition 1.1, citing reference [6], and to a further stable-category conjecture citing reference [12]. The corresponding bibliography entries and the further conjecture are not supplied.
\end{itemize}
\end{document}
